\documentclass{ximera} 

\input{../../preamble.tex}
% marges:
% \usepackage[paper=a4paper,margin=3.4cm,marginparwidth=2cm]{geometry}
% \usepackage[margin=.5in, includehead, includefoot, hmargin={.8in,.5in}, a4paper]{geometry}
% zes voorkeuren die je meteen zelf kan instellen:
\def\uitbr{0} % waarde 1 als je Uitbreiding in de marge wilt, 0 als je dat niet wilt 
\def\wsg{0} % waarde 1 als je verwijzingen naar Wiskunde Samen gevat² wilt, 0 als je dat niet wilt (onderstippelde woorden in de hoofdtekst en logo met paginanummer en hyperlink naar boekpagina in de marge)
\def\rectoverso{0} % waarde 1 als je het PDF-bestand recto-verso wilt laten afdrukken, 0 als je het recto wilt
\def\voetn{0} % waarde 1 als je de voetnoten wilt, 0 als je dat niet wilt 
\def\grm{0} % waarde 1 als je de schermafdrukken van TI-84 Plus wilt, 0 als je dat niet wilt
\def\grmlogo{0} % waarde 1 als je logo TI-84 Plus in de marge wilt, 0 als je dat niet wilt (met hyperlink naar filmpje) 



% packages algemeen:
\pdfOnly{
% \usepackage[english,dutch]{babel}
}
\usepackage{amsfonts}
\usepackage{amsmath} %  o.a. \eqref en \text, omgeving equation*
\usepackage{amssymb} % o.a. \nmid
\usepackage{amsthm} % o.a. omgeving proof
\usepackage{graphicx} % o.a. figuren
\usepackage{xcolor} % \color
% \usepackage[disable]{todonotes} % \todo 
\usepackage{comment} % omgeving comment
\usepackage{xcolor} % kleuren 
\usepackage[skip=7pt plus1pt, indent=0pt]{parskip} % skip = verticale ruimte tussen twee alinea's, indent = insprong bij nieuwe alinea
\usepackage{multicol} % kolommen
\usepackage{xfrac} % breuken

\usepackage{siunitx} % SI eenheden
\usepackage{eso-pic} % achtergrond voorpagina
\usepackage{tabularx} % kolommen voor rekenmachine
\usepackage{rotating} % voor omgeving turn

% % voor figuren met PSTricks:
% \usepackage{pstricks} 
% \usepackage{pstricks-add}
% \usepackage{pst-plot}
% \usepackage{pst-node}
% \usepackage{pst-coil}
% % \usepackage{auto-pst-pdf}

% voor figuren met TikZ:
\usepackage{tikz} 
\usepackage{tkz-euclide} 
\usepackage{pgfplots} 
\usetikzlibrary{calc,intersections,through,backgrounds,patterns} 
\pgfplotsset{compat=newest}
\usepgfplotslibrary{fillbetween,colormaps}
\usepackage{import}
\usepackage{tikz-3dplot}

% zelfde parskip in minipage:
\makeatletter
\setlength{\parskip}{\medskipamount}
\newcommand{\@minipagerestore}{\setlength{\parskip}{\medskipamount}}
\makeatother

% voor veranderen van de naam Bibliografie naar Referentielijst:
\pdfOnly{
\addto\captionsdutch{\renewcommand{\bibname}{Referentielijst}}

% voor trefwoordenlijst:
\usepackage{imakeidx}
\makeindex[title=Trefwoordenlijst,program=makeindex,options=-s index.tex,columns=2,intoc=true]
}
% voor nummeren van vergelijkingen:
%WIM%\numberwithin{equation}{chapter} % als je dit desactiveert dan worden vergelijkingen in bijvoorbeeld hoofdstuk 3 genummerd als (1), (2) etc. in plaats van (3.1), (3.2) etc.

% voor hyperlinks en bladwijzers in PDF-bestand:
%WIM%\usepackage[bookmarksopen,bookmarksopenlevel=0,hypertexnames=false,pdfa,bookmarksnumbered]{hyperref} 
%WIM%\usepackage{bookmark} 

% voor hyperlinks: 
\hypersetup{pdfborder={0 0 0}, pdfstartpage=1,linkbordercolor={1 0 0}}%pdfborder={0 0 0} is geen rand rond hyperlinks, pdfborder={1 1 1} wel

% korter commando voor displaystyle, om bijvoorbeeld breuken groter te drukken in doorlopende tekst: 
\newcommand{\D}{\displaystyle}

% (veel)gebruikte verzamelingen:
\newcommand\NN{\mathbb{N}} % verzameling van de natuurlijke getallen
\newcommand\QQ{\mathbb{Q}} % verzameling van de rationale getallen
\newcommand\RR{\mathbb{R}} % verzameling van de re\"ele getallen
\newcommand\ZZ{\mathbb{Z}} % verzameling van de gehele getallen

% (veel)gebruikte operatoren:
\def\co{\operatorname{co}} % coördinaat van een punt
\def\ggd{\operatorname{ggd}} % positieve grootste gemene deler van twee gehele getallen niet beide nul
\def\gr{\operatorname{gr}} % graad van een veelterm

% voor kleur grafieken (donkergroen):
\definecolor{graf}{RGB}{0,100,0} 

% voor lijsten: 
% \usepackage{enumerate}
% \usepackage[shortlabels]{enumitem} 
\setlist{topsep=0em, itemsep=-0.15em}

% voor small bullet:
\newcommand\sbullet[1][.5]{\mathbin{\vcenter{\hbox{\scalebox{#1}{$\bullet$}}}}}

% voor meer verticale ruimte tussen vergelijkingen in align
\addtolength{\jot}{0.1cm}

% voor meervoudige voetnoten:
\usepackage{fnpct}

% voor het onderdrukken van voetnoten als optie:
\usepackage{letltxmacro}
\LetLtxMacro\Oldfootnote\footnote
\newcommand{\EnableFootNotes}{%
  \LetLtxMacro\footnote\Oldfootnote%
}
\newcommand{\DisableFootNotes}{%
  \renewcommand{\footnote}[2][]{\relax}
}

% voor verticale lijn in de marge (uitbreiding):
\usepackage[framemethod=default]{mdframed}
\usepackage{marginnote}
\reversemarginpar
\ifthenelse{\uitbr < 1}{\definecolor{rood}{RGB}{255,255,255}}{\definecolor{rood}{RGB}{254,64,64}} %HEX: #fe4040
\mdfdefinestyle{uitbreiding}{%
    topline=false,
    rightline=false,
    bottomline=false,
    leftline=true,
    linecolor=rood,
    linewidth=5pt,
    rightmargin=0pt,
    skipabove=10pt,% ipv 3
    skipbelow=0pt,
    leftmargin=-25pt,
    innerleftmargin=20pt,
    innerrightmargin=0pt,
    innertopmargin=0pt,
    innerbottommargin=0pt%,
%	needspace=30pt %minimumhoogte vooraleer lijn wordt gesplitst
    }
% \newenvironment{Uitbreiding}{
%     \marginpar{
%         \center
% 		\vspace{0.1cm}
% 		\vspace{7pt}
% 		\rotatebox{90}{\color{rood}\Large \bf Uitbreiding}
% 	}
%     \begin{mdframed}[style=uitbreiding]
%     }{\vspace{-0.05cm}
%     \end{mdframed}
% }

% \newenvironment{Uitbreiding}{}{}
\ifdefined\HCode
\newenvironment{Uitbreiding}{%
    \begin{expandable}{xmuitweiding}{Uitbreiding:}%
}{%
    \end{expandable}%
}%
\else
\newtheorem{Uitbreiding}{Uitbreiding}
\renewcommand{\theUitbreiding}{} % titel instellen
\tcolorboxenvironment{Uitbreiding}{
    boxrule=0pt,
    boxsep=0pt,
    colback={gray!2!white},
    enhanced jigsaw,
    borderline west={2pt}{0pt}{red!55!black},
    sharp corners,
    before skip=10pt,
    after skip=10pt,
    breakable,
    % before upper={\vspace*{-0.45cm}},
}
\fi

        %\begin{xmuitweiding}}{\end{xmuitweiding}}
% omgevingen voor lemma, definitie, voorbeeld etc.:
\newtheoremstyle{mystyle}
    {0em} % Space above
    {0em} % Space below
    {\itshape} % Body font
    {} % Indent amount
    {\bfseries} % Theorem head font
    {.} % Punctuation after theorem head
    {.5em} % Space after theorem head
    {} % Theorem head spec (can be left empty, meaning `normal')
\theoremstyle{mystyle}
%WIM%\newtheorem{lemma}{Lemma}[chapter] % als je [chapter] desactiveert dan Voorbeeld 3 in plaats van Voorbeeld 5.3, als je [chapter] vervangt door [section] dan Voorbeeld 5.2.3 in plaats van Voorbeeld 5.3 
%\theoremstyle{definition} % als je dit activeert dan is wat in de omgeving staat niet cursief gedrukt
\newtheorem{oefening}[lemma]{Oefening}
\newtheorem{definitie}[lemma]{Definitie} 
\newtheorem{voorbeeld}[lemma]{Voorbeeld}
\newtheorem{eigenschap}[lemma]{Eigenschap} 
\newtheorem{stelling}[lemma]{Stelling} 
\newtheorem{gevolg}[lemma]{Gevolg} 
\newtheorem{afspraak}[lemma]{Afspraak} 
\newtheorem{werkwijze}[lemma]{Werkwijze}

% omgeving proof, aangepaste ruimte: 
\makeatletter
\renewenvironment{proof}[1][\proofname]{\par
  \vspace{-\topsep}% remove the space after the theorem
  \pushQED{\qed}%
  \normalfont
  \topsep0pt \partopsep0pt % no space before
  \trivlist
  \item[\hskip\labelsep
        \itshape
    #1\@addpunct{.}]\ignorespaces
}{%
  \popQED\endtrivlist\@endpefalse
  \addvspace{0pt plus 0pt} % no space after
}
\makeatother

% kaderstijlen uit SOHO Wiskunde Plantyn:
\colorlet{steunkleur}{black}
\colorlet{steunkleurlicht}{steunkleur!30!white}
\colorlet{steunkleurkader}{steunkleur!7!white}
\colorlet{steunkleurkaderlicht}{steunkleur!2!white}

\mdfdefinestyle{kaderstijl_vol_licht}
{skipabove=6pt, 
skipbelow=6pt, 
backgroundcolor=steunkleurkader,
linecolor=steunkleurlicht,
linewidth = 0.4pt, 
topline=true,
bottomline=true, 
rightline=true,
innerleftmargin=5pt,
innerrightmargin=5pt,
innertopmargin=5pt,
leftmargin=0cm,
rightmargin=0cm,
innerbottommargin=5pt,
needspace=30pt % minimumhoogte voor splitsen kader
}
\surroundwithmdframed[style=kaderstijl_vol_licht]{definitie}
\surroundwithmdframed[style=kaderstijl_vol_licht]{stelling}
\surroundwithmdframed[style=kaderstijl_vol_licht]{lemma}
\surroundwithmdframed[style=kaderstijl_vol_licht]{eigenschap}
\surroundwithmdframed[style=kaderstijl_vol_licht]{gevolg}

% voor omgevingen voor oefening en antwoord:
\newenvironment{Oefening}{%
    \begin{enumerate}[ 
    series=Oef,
    resume=Oef,
    leftmargin=1.78em,
    label={\bfseries\arabic*.},
    ref=\arabic*
    ]
    \item %
    }{%
    \end{enumerate}
}
\newenvironment{Antwoord}{%
    \begin{enumerate}[%
    series=Antw,
    resume=Antw,
    leftmargin=1.78em,
    label={\bfseries\arabic*.},
    ref=\arabic*
    ]
    \item %
    }{%
    \end{enumerate}
}

% als de omgeving Uitbreiding start met een kaderomgeving (definitie, stelling, eigenschap, lemma of gevolg) dan moet wat extra verticale ruimte voorzien worden, met het commando \uitbreidingstartmetkader:
\def\uitbreidingstartmetkader{\mbox{}\vspace{-0.205cm}}

%WIM % voor schema van de staartdeling:
\usepackage{stackengine}
\setstackgap{S}{5pt}
\stackMath\def\stackalignment{r}
\newcommand{\myRule}[3][white]{\textcolor{#1}{\rule{#2}{#3}}}
\let\ph\phantom % enkel voor tekst
\newcommand{\mph}[1]{% enkel voor math mode
    \mathcolor{white}{#1}%
}
\def\staartmin{\rule{0.25cm}{0.1mm}\myRule{0.3cm}{0.1mm}}
\def\staartphmin{\myRule{0.25cm}{0.1mm}\myRule{0.3cm}{0.1mm}}
\newcommand{\staartstreep}[1]{\rule{\widthof{$#1$}}{0.1mm}}
\newcommand{\staartphstreep}[1]{\myRule{\widthof{$#1$}}{0.1mm}}

% voor kolommen met schema van Horner:
\newcommand{\kolbreed}{1.0cm}
\newcolumntype{H}{>{\centering\arraybackslash$} p{\kolbreed} <{$}}

% voor nieuw commando utikzdashed: onderlijn (zoals underline) maar dan met stippellijn:
\tikzset{
    cheating dash/.code args={on #1 off #2 ends #3}{%
        \csname tikz@addoption\endcsname{%
            \pgfgetpath\currentpath%
            \pgfprocessround{\currentpath}{\currentpath}%
            \csname pgf@decorate@parsesoftpath\endcsname{\currentpath}{\currentpath}%
            \pgfmathparse{max(#1-#3,0)}\let\dashphase=\pgfmathresult%
            \pgfmathparse{\csname pgf@decorate@totalpathlength\endcsname-#1+2*\dashphase}\let\rest=\pgfmathresult%
            \pgfmathparse{#1+#2}\let\onoff=\pgfmathresult%
            \pgfmathparse{max(floor(\rest/\onoff), 1)}\let\nfullonoff=\pgfmathresult%
            \pgfmathparse{max((\rest-\onoff*\nfullonoff)/\nfullonoff+#2, #2)}\let\offexpand=\pgfmathresult%
            \pgfsetdash{{#1}{\offexpand}}{\dashphase pt}}%
    },
    cheating dash per segment/.style args={on #1 off #2 ends #3}{
        /utils/exec=\csname tikz@options\endcsname,%
        decoration={show path construction,
            lineto code={\draw [cheating dash=on #1 off #2 ends #3] (\tikzinputsegmentfirst) -- (\tikzinputsegmentlast);},
            curveto code={\draw [cheating dash=on #1 off #2 ends #3] (\tikzinputsegmentfirst) .. controls (\tikzinputsegmentsupporta) and (\tikzinputsegmentsupportb) .. (\tikzinputsegmentlast);},
            closepath code={\draw [cheating dash=on #1 off #2 ends #3] (\tikzinputsegmentfirst) -- (\tikzinputsegmentlast);}
        },
        decorate,
    },
}
\newcommand{\utikzdash}[1]{%
    \tikz[baseline=(todotted.base)]{
        \node[inner sep=0pt,outer sep=1.5pt] (todotted) {#1};
        \draw[cheating dash per segment=on 2pt off 2pt ends 2pt, line width=0.4pt] (todotted.south west) -- (todotted.south east);
    }%
}%

\providecommand{\xmemph}[1]{\textit{#1}}

% voor nieuw commando underdashed: onderlijn met stippellijnen en ook het woord cursief zetten:
\newcommand{\underdashed}[1]{%
    {\\textit{\utikzdash{\!#1\!}}}%
}

% voor icoon TI-84 Plus met verwijzing naar filmpje:
\newcommand{\grmlink}{\raisebox{0cm}{\includegraphics[width=1cm]{TI84Plus-icoon}}}
\newcommand{\grmref}[1]{%
	\ifthenelse{\grmlogo < 1}{}{
		\reversemarginpar%
		\marginpar{
			\vspace{-0.4cm}%
			\htmladdnormallink{\grmlink}{#1}
		}
	}	
}

% GRM knoppen:
\newcommand{\GRM}[1]{\fbox{\rule[0mm]{0cm}{0.215cm}\textup{\texttt{#1}}}}
\newcommand{\wedgetext}{{\raisebox{0.02cm}{\begin{turn}{90}>\end{turn}}}}
\newcommand{\veetext}{\raisebox{0.2cm}{\begin{turn}{-90}>\end{turn}}}

% voor kolommen met GRM screens:
\newlength{\widthallscreens}
\newlength{\widthscreens}
\newlength{\spaceleftscreen}
\newlength{\spacerightscreen}
\newlength{\spacebetweenscreens}

\newcommand{\setscreens}{
	\setlength{\spaceleftscreen}{2pt}
	\setlength{\spacerightscreen}{2pt}
	\setlength{\spacebetweenscreens}{8pt}
	\addtolength{\linewidth}{-28pt} % 3*8pt tussen vier screens en 2 pt links en 2 pt rechts
	\setlength{\widthallscreens}{\linewidth}
	\setlength{\widthscreens}{0.25\widthallscreens}
	\addtolength{\linewidth}{28pt}
	\newcolumntype{G}{p{\widthscreens}}
	\newcolumntype{s}{p{\spacebetweenscreens}}
	\newcolumntype{L}{p{\spaceleftscreen}}
	\newcolumntype{R}{p{\spacerightscreen}}
	\setlength{\tabcolsep}{0pt}
}

% voor icoon met verwijzing naar Wiskunde Samen gevat² in de marge:
\newcommand{\wsglink}{\raisebox{0cm}{\includegraphics[width=1cm]{wsglogo}}}
\newcommand{\htmladdnormallink}[2]{\href{#2}{#1}}
\newcounter{pagnrwsg}
\newcommand{\wsgref}[3]{% 
    % #1 woord dat onderstippeld wordt
    % #2 pagina van Wiskunde Samen gevat² waar je op terecht komt als je op de link klikt  
    % #3 afgedrukt op het logo van Wiskunde Samen gevat² in de marge: één of meerdere paginanummers
    \ifthenelse{\wsg < 1}{#1}{%
        \underdashed{#1}%
        {\setcounter{pagnrwsg}{#2}}%
        {\addtocounter{pagnrwsg}{14}}%
        \reversemarginpar%
        \marginpar{\vspace{-0.6cm}%
           \htmladdnormallink{\wsglink}{https://online.fliphtml5.com/sanky/laea/\#p=\arabic{pagnrwsg}}%
            \raisebox{0.41cm}[0cm][0cm]{%
                \hspace{-1.5cm}\makebox[2cm][c]{\colorbox{white}{\texttt{\footnotesize{#3}}}}%
            }%
        }%
    }%
}

% voor invoegen blanco pagina bij de optie recto-verso:
\def\blancobijrectoverso{
	\ifthenelse{\rectoverso < 1}{\clearpage}{
	\clearpage
	\thispagestyle{empty}
	\mbox{}
	\clearpage
	}
}

% aantal pagina's van het bestand:
\ifthenelse{\rectoverso < 1}{\def\totpag{57}}{\def\totpag{66}}

% documenteigenschappen:
% \usepackage{hyperxmp}
% \hypersetup{
% pdftitle={Open Source Wiskunde Aan zet: Veeltermen}, 
% pdfnumpages={\totpag},
% pdfauthor={Koen De Naeghel},
% pdflang={nl},
% pdfkeywords={wiskunde, open source, wiskunde aan zet, veeltermen, secundair onderwijs, tweede graad, doorstroomfinaliteit},
% pdfsubject={Veeltermen},
% pdfcopyright={\unichar{"24B8} 2024 Koen De Naeghel},
% pdfdate={17 december 2024},
% pdfapart=1,
% pdfstartview=}

% voor het benoemen van titel, auteur en datum:
% \title{Veeltermen}
% \author{Auteur: Koen De Naeghel}
% \date{\today}

\newcommand\BackgroundPic{
    \put(-260,-125){
    \parbox[b][\paperheight]{\paperwidth}{%
    \vfill
    \centering
    \includegraphics[height=\paperwidth, keepaspectratio]{WaZlogo}%
    \vfill
}}}


\newcommand{\chapter}{\section}

\addPrintStyle{..}

\begin{document}
	\author{Koen De Naeghel}
	\xmtitle{Veeltermen ontbinden in factoren}{}
	\label{xim:veeltermen_ontbinden_in_factoren}
%%\section{Veeltermen ontbinden in factoren}\index{veelterm!ontbinden in factoren}\index{ontbinden in factoren}

Tot nu toe heb je al enkele strategie\"en gezien om te \wsgref{ontbinden in factoren}{56}{56-60}: afzonderen, merkwaardige producten herkennen, groeperen en ontbinden van een kwadratische veelterm.

\begin{example}
De volgende kwadratische veelterm kan ontbonden worden in lineaire factoren omdat de discriminant positief is. Het rekenwerk kan ingekort worden door eerst een gemeenschappelijke factor van de termen af te zonderen. 
\begin{align*}
6x^2+21x-45 & = 3(2x^2 + 7x - 15) \\
& \qquad 
\begin{array}{|l}
\hline
\begin{aligned}
\vrule height 0.4cm width 0cm
& D = b^2 - 4ac = 7^2-4\cdot 2 \cdot (-15) = 169 > 0 \\[0.1cm] 
& x_{1} = \frac{-b + \sqrt{D}}{2a} = \frac{-7 + 13}{4} = \frac{3}{2} \\[0.1cm] 
& x_{2} = \frac{-b - \sqrt{D}}{2a} = \frac{-7 - 13}{4} = -5 \\[0.1cm]
\end{aligned} \\
\hline
\end{array} \\
& = 3 \cdot a(x-x_1)(x-x_2) \\
& = 3 \cdot 2\left(x- \frac{3}{2}\right)(x+5) \\
& = 3(2x-3)(x+5)
\end{align*}
\end{example}

In deze paragraaf leer je nog twee andere strategie\"en. Door al die werkwijzen doordacht in te zetten, kun je dan al heel wat eenvoudige veeltermen algebra\"isch ontbinden in factoren. Vooreerst vermelden we een belangrijk resultaat, dat een gevolg is van de zogenaamde \xmemph{hoofdstelling van de algebra}.\footnote{De hoofdstelling van de algebra stelt dat elke \xmemph{complexe} veelterm (een veelterm waarvan de co\"effici\"enten \xmemph{complexe getallen} zijn) dat niet constant is kan geschreven worden als een product van lineaire complexe veeltermen. Die stelling werd voor het eerst aangetoond door Carl Friedrich Gauss\index{Gauss, Carl Friedrich} in 1798.}\index{hoofdstelling van de algebra} Het bewijs laten we achterwege.

\begin{theorem}\label{stelling:ontbinden}
Elke veelterm kan geschreven worden als een product van constante veeltermen, lineaire veel\-termen en kwadratische veeltermen met strikt negatieve discriminant.
\end{theorem}

\begin{example}
Gegeven is telkens een veelterm met bijbehorende ontbinding in factoren. Om aan te geven dat we een kwadratische veelterm niet verder kunnen ontbinden in lineaire factoren, duiden we aan dat de discriminant $D$ strikt negatief is. Je kan elke gelijkheid controleren door het rechterlid uit te werken en te vereenvoudigen. 
\begin{enumerate}

\item
$\D 2x^3 + 5x^2 + 14x + 6 = (2x+1)\underbrace{(x^2+2x+6)}_{D < 0}$
\item
$\D x^4+x^2+1 = \underbrace{(x^2+x+1)}_{D < 0}\underbrace{(x^2-x+1)}_{D < 0}$
\end{enumerate}
\end{example}

%%% \clearpage

Sommige veeltermen kun je meteen ontbinden omdat je een merkwaardig product herkent. Zo volgt bijvoorbeeld uit $x^2 - a^2 = (x-a)(x+a)$ met $a \in \RR$ meteen dat
\[
x^2 - 7 = x^2 - \bigl(\sqrt{7}\bigr)^2 = \bigl(x-\sqrt{7})(x+\sqrt{7}\bigr). 
\]
Om ook een veelterm $A(x)$ van de vorm $x^3 - a^3$ te ontbinden in factoren, kunnen we als volgt te werk gaan.
We stellen vast dat $A(a) = 0$. Daarom is $A(x)$ deelbaar door $x-a$, zie Stelling \ref{stelling:kenmerkvandeelbaarheid} (kenmerk van deelbaarheid door $x-a$). We berekenen het quoti\"ent met het schema van Horner
\renewcommand{\kolbreed}{\widthof{$-a^3$}}
\[
\begin{array}{c|HHHH}
  & 1 & 0 & 0 & -a^3 \\[0.2cm]
\D a & \downarrow  & a  & a^2  & a^3  \\[0.2cm]
\hline 
\vrule height 1.2em width 0pt 
  & 1 & a & a^2 & \multicolumn{1}{|c}{0} 
\end{array}
\]
waaruit we afleiden dat $x^3 - a^3 = (x-a)(x^2 + ax + a^2)$. Zo kunnen we ook elk verschil van derde machten ontbinden in factoren, zoals
\[
x^3 - 125 = x^3 - 5^3 = (x-5)(x^2 + 5x + 25) 
\]
en ook
\[
8x^3 - 27 = (2x)^3 - 3^3 = (2x-3)(4x^2+6x+9).
\]
Op een gelijkaardige manier kun je ook een veelterm $A(x)$ van de vorm $x^3 + a^3$ ontbinden in factoren: deze keer is $A(-a) = 0$ zodat $A(x)$ is deelbaar door $x-(-a) = x+a$. Het bijbehorende schema van Horner is dan
\renewcommand{\kolbreed}{\widthof{$-a^3$}}
\[
\begin{array}{c|HHHH}
  & 1 & 0 & 0 & -a^3 \\[0.2cm]
\D -a & \downarrow  & -a  & a^2  & -a^3  \\[0.2cm]
\hline 
\vrule height 1.2em width 0pt 
  & 1 & -a & a^2 & \multicolumn{1}{|c}{0} 
\end{array}
\]
waaruit volgt dat  $x^3 + a^3 = (x+a)(x^2-ax+a^2)$. Zo is dan bijvoorbeeld
\[
x^3+8 = x^3 + 2^3 = (x+2)(x^2-2x+4).
\]
Bij uitbreiding kun je ook producten herkennen voor $x^4-a^4$ en $x^5 \pm a^5$. Die vatten we hieronder samen. Ook hier kun je elke gelijkheid aantonen door het rechterlid uit te werken en te vereenvoudigen. Je kan ook een gelijkheid terugvinden vanuit het linkerlid: pas het kenmerk van deelbaarheid door $x \pm a$ toe en voer het schema van Horner uit, zoals we hierboven voor $x^3 \pm a^3$ hebben laten zien.\footnote{Merk op dat we de producten voor $x^3 + a^3$ en $x^5 + a^5$ kunnen verkrijgen door in de gelijkheden voor $x^3 - a^3$ en $x^5 - a^5$ elke $a$ te vervangen door $-a$. Zoals Stelling \ref{stelling:ontbinden} hierboven aangeeft, kun je de tweetermen $x^4 - a^4$ en $x^5 \pm a^5$ nog verder ontbinden in factoren tot een product van lineaire en kwadratische veeltermen. Je kan ook $x^4+a^4$ ontbinden in factoren. Die ontbindingen staan vermeld in Oefening \ref{somtweevierdemachten}.} 
 
\begin{proposition}[merkwaardige producten herkennen]\index{merkwaardige producten}
Voor elke $a \in \RR$ is:
\begin{align*}
& x^2 - a^2 = (x-a)(x+a) && \text{verschil van twee kwadraten} \\
& x^3 + a^3 = (x+a)(x^2 - ax + a^2) && \text{som van twee derde machten} \\
& x^3 - a^3 = (x-a)(x^2 + ax + a^2) && \text{verschil van twee derde machten} \\
& x^4 - a^4 = (x-a)(x^3 + ax^2 + a^2x + a^3) && \text{verschil van twee vierde machten} \\
& x^5 + a^5 = (x+a)(x^4 - ax^3 + a^2x^2 - a^3x + a^4) && \text{som van twee vijfde machten} \\
& x^5 - a^5 = (x-a)(x^4 + ax^3 + a^2x^2 + a^3x + a^4) && \text{verschil van twee vijfde machten.}
\end{align*}
\end{proposition}

%%% \clearpage

Nu kunnen we sommige eenvoudige veeltermen zo ver mogelijk ontbinden in factoren door alle eerder geziene strategie\"en voor ontbinden in factoren te combineren (afzonderen, merkwaardige producten herkennen, groeperen en ontbinden van een kwadratische veelterm).

\begin{example}
We ontbinden telkens de veelterm zo ver mogelijk in factoren. 
\begin{enumerate}

\item 
$\D x^3 + 8 = x^3 + 2^3 = (x+2)\underbrace{(x^2-2x+4)}_{D = -12 < 0}$
\item
$\D -500x^3 + \frac{27}{2} = -\frac{1}{2}(1000x^3 - 27) = -\frac{1}{2}(10x-3)\underbrace{(100x^2 + 30x + 9)}_{D = -2700 < 0}$
\item
$\D x^4-25 = (x^2 - 5)(x^2+5) = (x-\sqrt{5})(x+\sqrt{5})\underbrace{(x^2+5)}_{D = -20 < 0}$
\item
$\D 9x^5 - 6x^4 - 16x^3 = x^3(9x^2-6x-16)$ 
\item[]
$\D \mph{9x^5 - 6x^4 - 16x^3} \qquad 
\begin{array}{|l}
\hline
\begin{aligned}
\vrule height 0.4cm width 0cm
& D = (-6)^2-4\cdot 9 \cdot (-16) = 612 = 2^2 \cdot 3^2 \cdot 17 \\[0.1cm] 
& x_{1,2} = \frac{-(-6) \pm \sqrt{6^2 \cdot 17}}{2 \cdot 9} = \frac{6 \pm 6\sqrt{17}}{18} = \frac{1 \pm \sqrt{17}}{3} \\[0.1cm] 
\end{aligned} \\
\hline
\end{array}
$
\item[]
$\D \mph{9x^5 - 6x^4 - 16x^3} = x^3 \cdot 9 \cdot \left(x-\frac{1+\sqrt{17}}{3}\right)\cdot\left(x-\frac{1-\sqrt{17}}{3}\right)$
\item[]
$\D \mph{9x^5 - 6x^4 - 16x^3} = x^3(3x-1-\sqrt{17})(3x-1+\sqrt{17})$
\end{enumerate}
\end{example}

Voor heel wat andere veeltermen volstaan die technieken niet, maar kan toch ontbonden worden door het kenmerk van deelbaarheid door $x-a$ toe te passen. Beschouw bijvoorbeeld
\[
A(x) = 2x^3 +x^2 - 8x + 21.
\]
Wegens het kenmerk van deelbaarheid (Stelling \ref{stelling:kenmerkvandeelbaarheid}) geldt voor elk re\"eel getal $a$: 
\[
(x-a) \mid A(x) \quad \Leftrightarrow \quad A(a) = 0. 
\]
Dus als $a$ een nulwaarde van de veelterm $A(x)$ is, dan kunnen we $A(x)$ ontbinden in factoren:
\[
2x^3 +x^2 - 8x + 21 = (x-a)(bx^2 + cx + d) % bx^3 + (c-ab)x^2 + (d-ac)x -ad => b = 2 geheel, c geheel, d geheel 
%\]
\quad 
\text{ voor zekere } b,c,d \in \RR.
\]
We merken op dat de co\"effici\"enten van de veelterm $A(x)$ gehele getallen zijn: $3$, $1$, $-8$ en $21$. In dat opzicht is het niet ondenkbaar dat de veelterm $A(x)$ een nulwaarde $a$ heeft dat zelf ook een geheel getal is. We spreken dan van een \xmemph{gehele nulwaarde}.

Als die veelterm $A(x)$ een gehele nulwaarde $a$ heeft, dan leert het uitwerken van het rechterlid ons dat de getallen $b,c,d$ zelf ook gehele getallen zijn. Bovendien vinden we dat $21 = -ad$, zodat $a$ een (gehele) deler is van $21$. Besluit: als $a$ een gehele nulwaarde van $A(x)$ is, dan is $a$ een deler van de constante term. 

Die redenering gaat ook op voor andere veeltermen met gehele co\"effici\"enten. Nulwaarden worden ook \xmemph{wortels} genoemd, en omdat het over gehele nulwaarden gaat, noemt men het resultaat de \xmemph{gehele wortelstelling}. Voor een bewijs verwijzen we naar Oefening \ref{oefening:gehelewortelstelling}.
\footnote{De voorwaarde $a \neq 0$ is nodig, omdat deelbaarheid door $0$ ongedefinieerd is.
Een veralgemening is de zogenaamde \xmemph{rationale wortelstelling}, zie Oefening \ref{oefening:rationalewortelstelling}.}

\begin{theorem}[gehele wortelstelling] \label{stelling:gehelewortelstelling}\index{gehele wortelstelling}
Zij $A(x)$ een veelterm met gehele co\"effici\"enten. Als $a \neq 0$ een gehele nulwaarde van $A(x)$ is, dan is $a$ een deler van de constante term. 
\end{theorem}

%%% \clearpage




De gehele wortelstelling leidt tot een techniek om elke veelterm met gehele co\"efici\"enten te ontbinden in factoren, op voorwaarde dat die veelterm een gehele nulwaarde heeft. Die nulwaarde(n) kun je opsporen door eerst de delers van de constante term te berekenen en daarna voor elke deler na te gaan of ze een nulwaarde is van de veelterm. In de praktijk kun je het berekenen van de getalwaarden uitvoeren met ICT. \footnote{De \xmemph{rationale wortelstelling} leidt tot een techniek om elke veelterm met \xmemph{rationale} co\"efici\"enten 
te ontbinden in factoren, op voorwaarde dat er een \xmemph{rationale} nulwaarde is. De kanshebbers voor rationale nulwaarden $\frac{p}{q}$ worden dan bepaald door de delers $p$ van de constante term en de delers $q$ van de hoogstegraadsco\"effici\"ent.}

\begin{example}
We ontbinden de volgende veelterm zo ver als mogelijk.
\renewcommand{\kolbreed}{\widthof{$-21$}}
\begin{align*}
A(x) & = 2x^3 + x^2 - 8x + 21 \\
& \qquad
\begin{array}{|l}
\hline
\vrule height 0.5cm width 0cm
\text{ kanshebbers gehele nulwaarden: delers van de constante term $21$
} \\[0.1cm]
\text{ ICT: } A(-3) = 0 \text{ dus $A(x)$ is deelbaar door $x- (-3) = x+3$} \\[0.1cm]
\text{ schema van Horner:} \\[0.1cm]
\qquad
\begin{array}{c|HHHH}
  & 2 & 1 & -8 & 21 \\[0.2cm]
\D -3 & \downarrow  & -6  & 15  & -21  \\[0.2cm]
\hline 
\vrule height 1.2em width 0pt 
  & 2 & -5 & 7 & \multicolumn{1}{|c}{0} 
\end{array} \\[-0.2cm]
\mbox{}\\
\hline
\end{array} \\[0.1cm]
& = (x+3)\underbrace{(2x^2-5x+7)}_{D = -31 < 0}
\end{align*}
\ifthenelse{\grm < 1}{}{
Gebruik van de grafische rekenmachine TI-84 Plus CE.
\grmref{https://www.dropbox.com/scl/fi/0gicz349b69zknu9fnj6o/TI84Plus-p.O-28.mp4?rlkey=e9x3cawyyw4zm6n3fbp33l6ie&raw=1}


\setscreens
\begin{tabularx}{\linewidth}{LGsGsGsGR}
& \GRM{y=} && \GRM{2nd}\,\GRM{table}\,\GRM{\veetext} && \GRM{\veetext} && \GRM{\wedgetext}\\[0.2cm]
&
{\includegraphics[width=\widthscreens]{screen01}}
&&
{\includegraphics[width=\widthscreens]{screen02}}
&&
{\includegraphics[width=\widthscreens]{screen03}}
&&
{\includegraphics[width=\widthscreens]{screen04}}
&
\end{tabularx}
\setlength{\tabcolsep}{6pt}
}
\end{example}

\begin{example} 
We ontbinden de volgende veelterm zo ver als mogelijk.
\renewcommand{\kolbreed}{\widthof{$-18$}}
\begin{align*}
A(x) & = x^3 + 4x^2 - 3x - 18 \\
& \qquad
\begin{array}{|l}
\hline
\vrule height 0.5cm width 0cm
\text{ kanshebbers gehele nulwaarden: delers van de constante term $-18$
} \\[0.1cm]
\text{ ICT: } A(2) = 0 \text{ en } A(-3) = 0 \text{ dus $A(x)$ is deelbaar door $(x-2)(x+3)$} \\[0.1cm]
\text{ schema's van Horner:} \\[0.1cm]
\qquad
\begin{array}{c|HHHH}
  & 1 & 4 & -3 & -18 \\[0.2cm]
2 & \downarrow  & 2  & 12  & 18  \\[0.2cm]
\hline 
\vrule height 1.2em width 0pt 
  & 1 & 6 & 9 & \multicolumn{1}{|c}{0} \\[0.2cm]
-3& \downarrow & -3 & -9 \\[0.2cm]
\cline{1-4}
\vrule height 1.2em width 0pt
  & 1 & 3 & \multicolumn{1}{|c}{0} 
\end{array} \\[-0.2cm]
\mbox{}\\
\hline
\end{array} \\[0.1cm]
& = (x-2)(x^2+6x+9) \\
& = (x-2)(x+3)(x+3) = (x-2)(x+3)^2
\end{align*}
\end{example}

%%% \clearpage

\begin{example} 
We ontbinden de volgende veelterm zo ver als mogelijk.
\renewcommand{\kolbreed}{\widthof{$-6$}}
\begin{align*}
& \mph{=} \frac{1}{2}\,x^4 + x^3 - x^2 - 3x - \frac{3}{2} \\
& = \frac{1}{2}\underbrace{\left(x^4 + 2x^3 - 2x^2 - 6x - 3 \right)}_{A(x)} \\
& \qquad
\begin{array}{|l}
\hline
\vrule height 0.5cm width 0cm
\text{ kanshebbers gehele nulwaarden: delers van de constante term $-3$
} \\[0.1cm]
\text{ ICT: } A(-1) = 0 \text{ dus $A(x)$ is deelbaar door $x+1$} \\[0.1cm]
\text{ schema van Horner:} \\[0.1cm]
\qquad
\begin{array}{c|HHHHH}
  & 1 & 2 & -2 & -6 & -3 \\[0.2cm]
\D -1 & \downarrow  & -1  & -1  & 3 & 3  \\[0.2cm]
\hline 
\vrule height 1.2em width 0pt 
  & 1 & 1 & -3 & -3 & \multicolumn{1}{|c}{0}
\end{array} \\[-0.2cm]
\mbox{}\\
\hline
\end{array} \\[0.1cm]
& = \frac{1}{2}(x+1)\underbrace{(x^3+x^2-3x-3)}_{Q(x)} \\
& \qquad
\begin{array}{|l}
\hline
\vrule height 0.5cm width 0cm
\text{ kanshebbers gehele nulwaarden: delers van de constante term $-3$
} \\[0.1cm]
\text{ ICT: } Q(-1) = 0 \text{ dus $Q(x)$ is deelbaar door $x+1$} \\[0.1cm]
\text{ schema van Horner:} \\[0.1cm]
\qquad
\begin{array}{c|HHHH}
  & 1 & 1 & -3 & -3 \\[0.2cm]
\D -1 & \downarrow  & -1  & 0  & 3  \\[0.2cm]
\hline 
\vrule height 1.2em width 0pt 
  & 1 & 0 & -3 & \multicolumn{1}{|c}{0}
\end{array} \\[-0.2cm]
\mbox{}\\
\hline
\end{array} \\[0.1cm]
& = \frac{1}{2}(x+1)(x+1)(x^2-3) \\
& = \frac{1}{2}(x+1)^2(x-\sqrt{3})(x+\sqrt{3})
\end{align*}
\end{example}
Vinden we slechts \'e\'en gehele nulwaarde $a$, dan kun je overwegen om  het schema van Horner meteen twee keer na elkaar uit te voeren. Is ook bij het tweede schema de rest gelijk aan nul, dan is de veelterm deelbaar door $(x-a)^2$. Bij het voorbeeld hierboven gaat dat als volgt.
\renewcommand{\kolbreed}{\widthof{$-6$}}
\begin{align*}
A(x) & = x^4 + 2x^3 - 2x^2 - 6x - 3 \\
& \qquad
\begin{array}{|l}
\hline
\vrule height 0.5cm width 0cm
\text{ kanshebbers gehele nulwaarden: delers van de constante term $-3$
} \\[0.1cm]
\text{ ICT: } A(-1) = 0 \text{ dus $A(x)$ is deelbaar door $x+1$} \\[0.1cm]
\text{ schema's van Horner:} \\[0.1cm]
\qquad
\begin{array}{c|HHHHH}
  & 1 & 2 & -2 & -6 & -3 \\[0.2cm]
\D -1 & \downarrow  & -1  & -1  & 3 & 3  \\[0.2cm]
\hline 
\vrule height 1.2em width 0pt 
  & 1 & 1 & -3 & -3 & \multicolumn{1}{|c}{0} \\[0.2cm]
-1 & \downarrow  & -1  & 0  & 3 \\[0.2cm]
\cline{1-5}
\vrule height 1.2em width 0pt 
  & 1 & 0 & -3 & \multicolumn{1}{|c}{0}
\end{array}\\[-0.2cm]
\mbox{}\\
\hline
\end{array} \\[0.1cm]
& = (x+1)(x+1)(x^2-3)
\end{align*}

%%% \clearpage
\end{document}